\chapter{Numerische Kriterien und Fehlertrennung}
\index{Konvergenz}\index{Residuum}\index{Bilanz}
\section{Normen und beobachtete Ordnung}
Für ein skalares Feld $w$ und seine diskrete Approximation $w_h$ seien
\begin{align}
 e_{L^2}&=\left(\int_\Omega |w-w_h|^2\,dV\right)^{1/2},\\
 e_{H^1}&=\left(\int_\Omega |\nabla w-\nabla w_h|^2\,dV\right)^{1/2}.
\end{align}
Die zweite Größe ist ein Halbnormfehler. Relative Fehler teilen durch die
entsprechende Norm der exakten Referenz; eine verschwindende Referenznorm
benötigt einen gesonderten absoluten Vergleich. Für Vektoren wird komponentenweise
summiert, für Gradienten die Frobeniusnorm verwendet. Ein Knotenmaximum ist
eine andere Zielgröße und darf nicht als integrierte Norm bezeichnet werden.

Für zwei aufeinanderfolgende Netze ergibt sich
\begin{equation}
 p_{h,H}=\frac{\log(e_H/e_h)}{\log(H/h)}.
\end{equation}
Der Nenner muss zum tatsächlichen Verfeinerungsfaktor passen. Die letzten zwei
Intervalle einer vorab gewählten Viernetzfolge dienen in den glatten räumlichen
Tests zur Prüfung der asymptotischen Ordnung. Positivität und monotone
Fehlerabnahme werden zusätzlich geprüft; nachträglich ausgewählte günstige
Intervalle wären kein gleichwertiger Nachweis.

\section{Räumliche Konvergenz}
Die glatten Tet4-Referenzen erwarten $L^2$-Ordnung 2 und $H^1$-Ordnung 1.
Typische vorab festgelegte Fenster sind $[1.8,2.2]$ und $[0.85,1.15]$.
Materialgrenzen und gekoppelte Tests besitzen ihre eigenen dokumentierten
Fenster; beispielsweise nutzt der historische Thermopiezo-Modelltest
$[1.75,2.25]$, $[0.75,1.25]$ und für den Energiefehler $[1.5,2.5]$.
Diese Kriterien werden nicht auf einen anderen Fall übertragen oder gelockert.

Tet10 nutzt in den glatten Feldtests vier Netze mit $N=1,2,3,4$ und die Fenster
$[2.7,3.3]$ für $L^2$ sowie $[1.8,2.2]$ für $H^1$. Der historische
M7-Abnahmebericht \path{Doc/55_TET10_ACCEPTANCE.md} dokumentiert am
28. September 2026 elf Konvergenztests. Sein damaliger Commit und seine damalige
Suite sind historische Nachweise, keine aktuelle Gesamtzahl.

\begin{table}[htbp]\centering
\begin{tabular}{lrr}\toprule
Historischer letzter Schritt $N=3\to4$ & $p_{L^2}$ & $p_{H^1}$\\\midrule
Wärme / Elektrostatik & 2.985 & 1.978\\
Elastizität / thermomechanisches $u$ & 2.989 & 1.980\\
Piezoelektrisches $u$ & 3.007 & 1.994\\
Thermopiezoelektrisches $u$ & 3.008 & 1.994\\
Piezoelektrisches $\phi$ / thermisches $T$ & 2.985 & 1.978\\
\bottomrule\end{tabular}
\caption{Gerundete Werte des versionierten M7-Berichts; glatte synthetische Felder.}
\end{table}
Sie sind keine Fehlergrenzen für reale Materialübergänge, Singularitäten an
Einspannungen oder beliebige unstrukturierte Netze.

\section{Zeitfehler und dynamische Zielgrößen}
Backward Euler besitzt hier Zeitordnung 1. Ein Tet10-Test mit exakt räumlich
darstellbarem $T=t^2$ trennt sie vom Raumfehler: bei 10, 20, 40 und 80 Schritten
ist der Endzeitfehler proportional zu $1/n$, mit Ordnungsfenster $[0.99,1.01]$.
Der Tet4-Abkühlungsfall verwendet das eigene Fenster $[0.9,1.1]$.
Kapazitäts- und Leistungsbilanz sowie zeitabhängige Randereignisse werden
gesondert geprüft. Ein Zeitkonvergenztest mit festem grobem Netz kann sonst
auf einem räumlichen Fehlerplateau enden.

Newmark mit $\beta=1/4$, $\gamma=1/2$ wird am analytischen Ein-DOF-Oszillator
geprüft. Bei Halbierung des Zeitschritts soll das Verhältnis der Fehler im
vorgegebenen Test zwischen 3.8 und 4.2 liegen, entsprechend Ordnung 2.
Energie und dynamisches Residuum sind zusätzliche Kriterien, keine alternative
Definition derselben Ordnung. Bewegte Dirichlet-Ränder und Lastwechsel
brauchen die korrekte Anfangsbeschleunigung und die konsistenten Randterme.

Tet10-Moden und harmonische Stabantworten zeigen für bestimmte skalare
Fehler der Frequenzen oder Endverschiebungen höhere Ordnungen. Das vorgegebene Fenster
$[3.5,4.5]$ beschreibt diese Zielgrößen und darf nicht als $H^1$-Ordnung 4
des gesamten Feldes ausgegeben werden. Nullfrequenz, Phasenlage, Ladungsbilanz
und Resonanzablehnung bleiben eigenständige Prüfungen.

\section{Residuen, Reaktionen und physikalische Bilanzen}
Für das freie lineare System lautet ein normierter Rückwärtsfehler
\begin{equation}
 \eta=\frac{\|Kx-f\|}{\|K\|\,\|x\|+\|f\|}.
\end{equation}
Ein kleines $\eta$ bedeutet, dass ein nahes algebraisches Problem gelöst wurde.
Die Konditionierung bestimmt, wie daraus ein Vorwärtsfehler entsteht.
Reaktionen werden aus dem ursprünglichen System vor Dirichlet-Elimination
berechnet. Bei skalierten gekoppelten Systemen muss für physikalische Residuen
auf die ursprünglichen Felder und Einheiten zurückgegangen werden.

Wärmebilanzen vergleichen Leistungen in W, Elektrostatik Ladungen in C und
Mechanik Kräfte in N bzw. Momente in N\,m. Energie in J ist eine andere Größe.
Ein gemeinsamer absoluter Grenzwert für gemischte Feldkomponenten hätte keine
einheitliche physikalische Bedeutung. Das Vorzeichen richtet sich nach dem
im Theorieband festgelegten äußeren Fluss und den Originalreaktionen.

Die Toleranzwahl berücksichtigt analytische Genauigkeit, Quadratur,
Rundung in Gleitkommaarithmetik, Größenordnung und Konditionierung. Ein Einheitsproblem mit
O(1)-Matrizen kann $10^{-12}$ tragen; starke Materialkontraste und gekoppelte
Blöcke benötigen eine eigene Skalierungsbetrachtung. Die historischen
Benchmarkbilanzen verlangen $10^{-9}$ in ihrer jeweiligen SI-Einheit.
Das ist weder eine universelle Solvervorgabe noch ein Netzfehlerbudget.

\section{Surrogatfehler und stationäre Studien}
\index{Surrogat!Nachweis}
Die M8-Studientests sichern Kennwertdefinitionen, Einheiten, Modell- und
Tabellenhashes, Fehlerstatus und Wiederaufnahme. Sie belegen nicht automatisch
die Netzunabhängigkeit jeder in einer Studie erzeugten Antwort.

Ein Surrogat benötigt getrennte Trainings- und Testfälle. Gleiche
Parametergruppen dürfen nicht über beide Mengen verteilt werden; eine
Trainingsgüte ist kein unabhängiger Testfehler. Die bestehenden Tests prüfen
analytische Interaktionen und einen echten FEM-Durchlauf mit sechs Trainings-
und zwei zurückgehaltenen Fällen. Ein absichtlich zu schwaches lineares Modell
für eine quadratische Antwort muss die Abweichung sichtbar berichten.

MAE, RMSE und maximaler absoluter Fehler besitzen die Einheit der Zielgröße.
Sie beschreiben den gewählten Testbestand. Ein Fehlerbalken für alle zukünftigen
Parameter folgt daraus nicht. Die zulässige Parameterbox garantiert außerdem
keine Abdeckung ihres gesamten Inneren durch Trainingspunkte.
Rangdefizite, ungültige Zeilen und Extrapolation werden ausdrücklich behandelt.
Surrogatfehler, FEM-Diskretisierungsfehler und experimentelle Modellabweichung
bleiben drei unterschiedliche Nachweise.
